
\documentclass{article}
\usepackage{colortbl}
\usepackage{makecell}
\usepackage{multirow}
\usepackage{supertabular}

\begin{document}

\newcounter{utterance}

\centering \large Interaction Transcript for game `adventuregame', experiment `new{-}words\_home{-}delivery\_easy\_basic\_easy', episode 5 with reka{-}flash{-}3.1{-}t0.0.
\vspace{24pt}

{ \footnotesize  \setcounter{utterance}{1}
\setlength{\tabcolsep}{0pt}
\begin{supertabular}{c@{$\;$}|p{.15\linewidth}@{}p{.15\linewidth}p{.15\linewidth}p{.15\linewidth}p{.15\linewidth}p{.15\linewidth}}
    \# & \multicolumn{2}{c}{Player} && \multicolumn{2}{c}{Game Master} \\
    \hline

    \theutterance \stepcounter{utterance}  
    & & & \multicolumn{4}{p{0.6\linewidth}}{
        \cellcolor[rgb]{0.9,0.9,0.9}{
            \makecell[{{p{\linewidth}}}]{
                \texttt{\tiny{[P1$\langle$GM]}}
                \texttt{You are playing a text adventure game. I will describe what you can perceive in the game. You write the single action you want to take in the game starting with >. Only reply with an action.} \\
\texttt{For example:} \\
\texttt{> open cupboard} \\
\\ 
\texttt{In addition to common actions, you can perce. To perce is to make something physically unobstructed and uncovered, and make its contents available for use or interaction.} \\
\\ 
\texttt{Your goal for this game is: Put the uner on the trion, the natin on the trion and the peach on the trion.} \\
\\ 
\texttt{Once you have achieved your goal, write "> done" to end the game.} \\
\\ 
\texttt{You are in a bedroom now. There are a pillow, a wardrobe, a bed and a book. The pillow is on the bed. The wardrobe is closed. The book is on the floor. There is a passage to a hallway here.} \\
            }
        }
    }
     \\ \\

    \theutterance \stepcounter{utterance}  
    & \multicolumn{4}{p{0.6\linewidth}}{
        \cellcolor[rgb]{0.561,0.737,0.561}{
            \makecell[{{p{\linewidth}}}]{
                \texttt{\tiny{[P1$\rangle$GM]}}
                \texttt{<reasoning>} \\
\texttt{Let's see. The goal is to put the "uner", "natin", and "peach" on the "trion". First, I need to figure out what these items are. The description mentions a pillow, wardrobe, bed, book, and a passage to a hallway. Maybe "uner" and "natin" are typos or misheard words? Let me check the original problem again.} \\
\\ 
\texttt{Wait, the user wrote: "Put the uner on the trion, the natin on the trion and the peach on the trion." Hmm, maybe "uner" and "natin" are objects in the game that I need to find. But the current room has a pillow, wardrobe, bed, book, and passage. The peach isn't mentioned here. Maybe the peach is in another room, like the hallway? Or perhaps "peach" is a typo for "peach" as in the fruit, but where would that be?} \\
\\ 
\texttt{Wait, the initial description says: "You are in a bedroom now. There are a pillow, a wardrobe, a bed and a book. The pillow is on the bed. The wardrobe is closed. The book is on the floor. There is a passage to a hallway here." So the items here are pillow, wardrobe, bed, book. No peach. So maybe the peach is in another room? The passage leads to a hallway. So perhaps I need to go to the hallway to find the peach. But the goal is to put all three items (uner, natin, peach) on the trion. Wait, but where is the trion? The current room has a bed, wardrobe, pillow, book. No trion mentioned. Maybe "trion" is another object that I need to find? Or perhaps "trion" is a typo for "tray" or something else?} \\
\\ 
\texttt{Wait, maybe "trion" is a specific object in the game. Let me check the problem statement again. The user says: "Your goal for this game is: Put the uner on the trion, the natin on the trion and the peach on the trion." So "trion" must be an object in the game. But in the current room description, there's no trion. The items are pillow, wardrobe, bed, book. So perhaps the trion is in another room, like the hallway? Or maybe the tr} \\
            }
        }
    }
    & & \\ \\

    \theutterance \stepcounter{utterance}  
    & & & \multicolumn{2}{p{0.3\linewidth}}{
        \cellcolor[rgb]{0.9,0.9,0.9}{
            \makecell[{{p{\linewidth}}}]{
                \texttt{\tiny{[GM$|$GM]}}
                \texttt{Round: 0} \\
            }
        }
    }
    & & \\ \\

    \theutterance \stepcounter{utterance}  
    & & & \multicolumn{2}{p{0.3\linewidth}}{
        \cellcolor[rgb]{0.9,0.9,0.9}{
            \makecell[{{p{\linewidth}}}]{
                \texttt{\tiny{[GM$|$GM]}}
                \texttt{command\_tag\_missing} \\
            }
        }
    }
    & & \\ \\

    \theutterance \stepcounter{utterance}  
    & & & \multicolumn{2}{p{0.3\linewidth}}{
        \cellcolor[rgb]{0.9,0.9,0.9}{
            \makecell[{{p{\linewidth}}}]{
                \texttt{\tiny{[GM$|$GM]}}
                {'goal_states_achieved': [], 'game_successfully_finished': False}
            }
        }
    }
    & & \\ \\

    \theutterance \stepcounter{utterance}  
    & & & \multicolumn{2}{p{0.3\linewidth}}{
        \cellcolor[rgb]{0.9,0.9,0.9}{
            \makecell[{{p{\linewidth}}}]{
                \texttt{\tiny{[GM$|$GM]}}
                \texttt{command\_tag\_missing} \\
            }
        }
    }
    & & \\ \\

\end{supertabular}
}

\end{document}
